%% kapa: exam=1 round=37 year=2026 period=2 form=A subject=accounting subject_ko=회계학 seq=23 number=63
%% source: Q-Net 감정평가사 기출문제 (cst00309) artlSeq=5258924 file=raw/1008/5258924/2026년 제37회 감정평가사 1차 2교시.pdf
%% conversion: kordoc 4.24.51 -> tools/gen1.py (auto); review=proofread-sonnet; flags=none
\begin{question}{63}
외화자산 및 외화부채의 환산방법에 관한 설명으로 옳지 않은 것은?
\begin{choices}
  \item 기능통화로 외화거래를 최초로 인식하는 경우에 거래일의 외화와 기능통화 사이의 현물환율을 외화금액에 적용하여 기록한다.
  \item 공정가치로 측정하는 비화폐성 외화항목은 공정가치가 측정된 날의 환율로 환산한다.
  \item 기능통화가 변경되는 경우에는 새로운 기능통화에 의한 환산절차를 변경한 날부터 전진적용한다.
  \item 여러 가지 환율을 사용할 수 있는 경우에는 해당 거래나 잔액에 따른 미래현금흐름이 측정일에 발생하였다면 결제하였을 환율을 사용한다.
  \item 역사적원가로 측정하는 비화폐성 외화항목은 마감환율로 환산하여 그 차이를 당기손익으로 인식한다.
\end{choices}
\end{question}
\answer{5}
